\documentclass[11pt,a4paper]{article}
\usepackage[utf8]{inputenc}
\usepackage[T1]{fontenc}
\usepackage{amsmath,amssymb}
\usepackage{siunitx}
\usepackage{booktabs}
\usepackage{geometry}
\usepackage{hyperref}
\geometry{margin=2.5cm}

\title{Single-Phase Power Transformer Sizing\\
\large A Practical Methodology for High-Frequency, High-Power Designs}
\author{}
\date{}

\begin{document}
\maketitle

\section{Scope}
This document summarizes the engineering steps required to dimension a single-phase
power transformer for switch-mode converter applications (e.g.\ PSFB, DAB, isolated
DC/DC stages), covering:
\begin{itemize}
    \item turns ratio determination from the voltage transfer function,
    \item core material selection as a function of operating frequency,
    \item core sizing via the area-product method,
    \item winding design considerations (skin/proximity effects, litz wire vs.\ foil),
    \item a worked numerical example at \SI{800}{V}/\SI{800}{V}, \SI{10}{kHz}, \SI{250}{kW}.
\end{itemize}

\section{Turns Ratio}
For an ideal isolated transformer driven by a quasi-square or square voltage waveform
(as typical of PSFB, DAB, or any hard/soft-switched bridge), the volt-second balance on
each winding gives Faraday's law in its waveform-corrected form:
\begin{equation}
    V = K_f \, N \, A_c \, B_{m} \, f
    \label{eq:faraday}
\end{equation}
where:
\begin{itemize}
    \item $V$ = RMS (or peak, depending on $K_f$ convention) voltage applied to the winding [V],
    \item $N$ = number of turns,
    \item $A_c$ = core cross-sectional area [m$^2$],
    \item $B_m$ = peak flux density [T],
    \item $f$ = switching/operating frequency [Hz],
    \item $K_f$ = waveform factor: $K_f = 4$ for a symmetric square wave, $K_f = 4.44$ for a pure sinusoid, intermediate values for trapezoidal/quasi-square waveforms with duty cycle $D<0.5$ (e.g.\ $K_f \approx 4D$ for a unipolar square wave of duty $D$).
\end{itemize}

The turns ratio is fixed by the required voltage transformation:
\begin{equation}
    n = \frac{N_1}{N_2} = \frac{V_1}{V_2}
\end{equation}
For converters requiring voltage step change under varying duty cycle (e.g.\ PSFB with
output inductor), $n$ is chosen at the nominal/maximum duty point, with margin for
line and load regulation.

\section{Core Material Selection vs.\ Frequency}
The dominant constraint at a given frequency is core loss, which is approximately:
\begin{equation}
    P_{core} \approx k_c \, f^{\alpha} \, B_m^{\beta} \, V_{core}
\end{equation}
(Steinmetz equation), where $\alpha,\beta$ and $k_c$ are material-dependent. The choice
of material is therefore frequency-driven:

\begin{center}
\begin{tabular}{@{}lllp{5.5cm}@{}}
\toprule
Material & Typical $f$ range & $B_{sat}$ [T] & Notes \\
\midrule
Grain-oriented Si-Fe (laminated) & DC--\SI{1}{kHz} & 1.6--2.0 & High saturation, but eddy-current losses scale with $f^2 t^2$; lamination thickness limits practical use above $\sim$\SI{1}{kHz}. \\
Amorphous (Fe-based) & \SI{1}{}--\SI{10}{kHz} & 1.5--1.6 & Lower loss than Si-Fe at the same $f$, good compromise for medium-frequency, high-power transformers. \\
Nanocrystalline (e.g.\ Vitroperm, Finemet) & \SI{1}{}--\SI{150}{kHz} & 1.2--1.5 & High $B_{sat}$ retained at MF, low specific losses; preferred for high-power MF transformers (DAB, solid-state transformers, traction). \\
MnZn ferrite & \SI{20}{kHz}--\SI{500}{kHz} & 0.35--0.5 & Low $B_{sat}$ limits power density at low kHz; becomes attractive above $\sim$\SI{50}{kHz} where its losses remain low and nanocrystalline/amorphous losses rise. \\
NiZn ferrite & $>$\SI{500}{kHz} & 0.3--0.4 & High-frequency, low-power applications; not relevant for high-power MF transformers. \\
\bottomrule
\end{tabular}
\end{center}

\textbf{Rule of thumb:} for a given frequency, select the material that allows the
highest usable $B_m$ (limited by core loss density, typically targeting
\SIrange{0.5}{1.5}{W/cm^3} depending on cooling) while satisfying thermal constraints.
At \SI{10}{kHz} and high power, nanocrystalline or amorphous tape-wound cores
typically outperform ferrite because their higher $B_{sat}$ allows a smaller, lighter
core for the same volt-seconds, even though their specific core losses per unit volume
at \SI{10}{kHz} are higher than ferrite's — the reduction in required volume usually
wins for power levels in the hundreds of kW.

\section{Core Sizing: Area-Product Method}
The area product $A_p = A_c \cdot A_w$ (core area times window area) relates the
apparent power transferred to the achievable current density and flux density,
and is the standard first-pass sizing tool (McLyman):
\begin{equation}
    A_p \;[\text{cm}^4] = \frac{P_t \cdot 10^4}{K_f \, K_u \, B_m \, f \, J}
    \label{eq:ap}
\end{equation}
where:
\begin{itemize}
    \item $P_t$ = total apparent power handled by the transformer, $P_t \approx P_{in}+P_{out} \approx 2P$ for an isolated transformer carrying power $P$ on each side,
    \item $K_u$ = window utilization factor ($\approx 0.2$--$0.3$ for round/litz wire with insulation and creepage, up to $\approx 0.4$--$0.5$ for foil/busbar windings),
    \item $B_m$ = peak flux density [T] (material- and loss-limited, see above),
    \item $f$ = frequency [Hz],
    \item $J$ = RMS current density [A/cm$^2$] (cooling-dependent: $\sim$200--350 for natural convection, up to 500+ with forced air/liquid cooling).
\end{itemize}
Once $A_p$ is known, a core (or stack of cores) is selected from a manufacturer
catalogue with matching or slightly larger $A_p$, giving $A_c$ and $A_w$ directly.
Equation~\eqref{eq:faraday} then yields the turns count for each winding.

\section{Winding Design}
At MF/HF, conduction losses are dominated by skin and proximity effects rather
than DC resistance:
\begin{equation}
    \delta = \sqrt{\frac{\rho}{\pi f \mu_0 \mu_r}}
\end{equation}
For copper, $\delta \approx 7.5/\sqrt{f\,[\text{Hz}]}$ cm. At \SI{10}{kHz},
$\delta \approx \SI{0.66}{mm}$, meaning conductor thickness should not exceed
roughly $2\delta \approx \SI{1.3}{mm}$ in any single layer.

\begin{itemize}
    \item \textbf{Litz wire}: appropriate for moderate current windings where many
    strands keep each strand below $\delta$; strand count and bundle geometry must
    be checked for proximity losses in multi-layer windings.
    \item \textbf{Copper foil}: preferred at very high current, low turn-count
    windings (as in the example below); foil thickness $\le 2\delta$, interleaved
    with thin insulation, often the only practical option above a few hundred
    amperes RMS per turn.
    \item \textbf{Interleaving}: primary/secondary interleaving reduces the MMF
    excursion across the window and is essential to control proximity losses at
    MF in high-current designs.
\end{itemize}

\section{Leakage Inductance Estimation}
Leakage inductance arises from the imperfect coupling between primary and
secondary windings and is set almost entirely by winding geometry rather than
core material. For concentric (or foil) windings arranged radially in the
window, the standard analytical estimate (Erickson \& Maksimovic;
McLyman)\cite{erickson,mclyman} referred to the primary is:
\begin{equation}
    L_{lk} \approx \frac{\mu_0 \, N_1^2 \, \mathrm{MLT}}{h_w}
    \left[ \frac{b_1+b_2}{3} + \delta \right]
    \label{eq:lk}
\end{equation}
where:
\begin{itemize}
    \item $\mathrm{MLT}$ = mean length per turn [m],
    \item $h_w$ = winding (window) height available for the turns [m],
    \item $b_1, b_2$ = total radial (or axial, depending on construction)
    build-up of the primary and secondary windings respectively [m],
    \item $\delta$ = thickness of the inter-winding insulation/isolation gap [m].
\end{itemize}
Equation~\eqref{eq:lk} shows the two key levers for leakage reduction:
\begin{itemize}
    \item \textbf{Increasing $h_w$} (taller, narrower window) spreads the same
    turns over a larger MMF-cancellation path, reducing $L_{lk}$.
    \item \textbf{Interleaving} the primary and secondary into $m$ alternating
    sections divides the effective leakage by approximately $m^2$:
    \begin{equation}
        L_{lk,\,interleaved} \approx \frac{L_{lk}}{m^2}
    \end{equation}
\end{itemize}
For DAB or PSFB topologies, leakage inductance is often a wanted design
parameter (it forms part or all of the power-transfer/resonant inductance), so
this calculation should be checked against the value required by the converter
control law, not only minimized.

\section{Insulation and Isolation}
Creepage and clearance distances are dictated by the working voltage, altitude,
pollution degree, and required isolation class (basic/supplementary/reinforced)
per IEC~60664-1 or the relevant converter standard. Tape-wound cores with bobbin
or potted construction are common; for busbar/foil windings, solid insulation
(e.g.\ polyimide, Nomex, epoxy-glass) must be verified for partial-discharge-free
operation at the working dV/dt, which can be severe with hard-switched square waves.

\section{Worked Example: \SI{800}{V}/\SI{800}{V}, \SI{10}{kHz}, \SI{250}{kW}}

\subsection*{Assumptions}
\begin{itemize}
    \item Square-wave excitation (PSFB/DAB-type bridge), $K_f = 4$.
    \item Turns ratio $n=1$ (\SI{800}{V}/\SI{800}{V}).
    \item Nanocrystalline tape-wound core, $B_m = \SI{1.0}{T}$ (derated from
    $B_{sat}\approx\SI{1.2}{T}$ for MF loss margin).
    \item $K_u = 0.3$ (foil winding, generous insulation/creepage at this power level).
    \item $J = \SI{300}{A/cm^2}$ (forced-air or liquid-cooled busbar/foil winding).
\end{itemize}

\subsection*{Step 1 — Apparent power}
\begin{equation}
    P_t = 2P = 2 \times \SI{250}{kW} = \SI{500}{kW}
\end{equation}

\subsection*{Step 2 — Area product}
\begin{equation}
    A_p = \frac{500{,}000 \times 10^4}{4 \times 0.3 \times 1.0 \times 10{,}000 \times 300}
    \approx \SI{1389}{cm^4}
\end{equation}
This is a large core (or, in practice, a paralleled/stacked assembly of standard
nanocrystalline C-cores or cut cores), consistent with a \SI{250}{kW} MF transformer.

\subsection*{Step 3 — Core area and turns}
Assuming a roughly square window/core split, $A_c \approx \sqrt{A_p} \approx \SI{37}{cm^2}$.
From Eq.~\eqref{eq:faraday}:
\begin{equation}
    N_1 = \frac{V_1}{K_f \, A_c \, B_m \, f}
    = \frac{800}{4 \times 37\times10^{-4} \times 1.0 \times 10{,}000}
    \approx 5.4 \;\Rightarrow\; N_1 = N_2 = 6 \ \text{turns}
\end{equation}
(rounding up gives margin against saturation; with $n=1$, $N_2=N_1$).

\subsection*{Step 4 — Conductor sizing}
RMS current per winding (lossless approximation):
\begin{equation}
    I_1 = I_2 = \frac{P}{V} = \frac{250{,}000}{800} = \SI{312.5}{A}
\end{equation}
Required copper cross-section at $J=\SI{300}{A/cm^2}$:
\begin{equation}
    A_{Cu} = \frac{I}{J} = \frac{312.5}{300} \approx \SI{1.04}{cm^2}
\end{equation}
With skin depth $\delta \approx \SI{0.66}{mm}$ at \SI{10}{kHz}, this area should be
realized as a foil (or paralleled foils) of thickness $\le \SI{1.3}{mm}$ and
sufficient width, or an equivalent litz bundle — foil/busbar is the practical
choice at this current level and turn count.

\subsection*{Step 5 — Leakage inductance estimate}
Assume a square-ish window of side $\approx\SI{6.1}{cm}$, so $h_w \approx \SI{6}{cm}$,
and a mean turn length $\mathrm{MLT}\approx\SI{20}{cm}$ (typical for this core size).
Each winding uses 6 foil turns of thickness $\SI{1.3}{mm}$ with $\SI{0.1}{mm}$
inter-turn insulation, giving a winding build-up:
\begin{equation}
    b_1 = b_2 = 6\times(1.3+0.1)\,\text{mm} = \SI{8.4}{mm}
\end{equation}
With an inter-winding isolation gap $\delta = \SI{2}{mm}$, Eq.~\eqref{eq:lk} gives:
\begin{equation}
    L_{lk} \approx \frac{4\pi\times10^{-7} \times 6^2 \times 0.20}{0.06}
    \left[\frac{0.0084+0.0084}{3}+0.002\right] \approx \SI{1.14}{\micro H}
\end{equation}
If the primary and secondary are interleaved into two sections each ($m=2$),
this drops to $L_{lk}\approx\SI{0.29}{\micro H}$. This value should be compared
against the leakage/resonant inductance required by the DAB/PSFB power-transfer
equation; if higher inductance is needed, an air gap or reduced interleaving can
be used deliberately to tune $L_{lk}$ upward.

\subsection*{Summary Table}
\begin{center}
\begin{tabular}{@{}ll@{}}
\toprule
Parameter & Value \\
\midrule
Core material & Nanocrystalline (e.g.\ Vitroperm 500F) \\
$B_m$ & \SI{1.0}{T} \\
Area product $A_p$ & $\approx \SI{1389}{cm^4}$ \\
$N_1 = N_2$ & 6 turns \\
$I_1 = I_2$ (RMS) & \SI{312.5}{A} \\
Winding type & Interleaved copper foil, $t \le \SI{1.3}{mm}$ \\
$L_{lk}$ (non-interleaved / $m{=}2$) & $\approx\SI{1.14}{\micro H}$ / $\SI{0.29}{\micro H}$ \\
\bottomrule
\end{tabular}
\end{center}

\subsection*{Notes}
This first-pass sizing must be followed by: (i) FEA cross-check of the analytical
leakage inductance estimate above (relevant for DAB/PSFB soft-switching range),
(ii) thermal verification of core and winding losses against the chosen cooling
method, and (iii) insulation coordination for the working voltage and dV/dt. The
area-product method gives the correct order of magnitude but final core selection
should match an actual manufacturer part (or stack thereof) with $A_p \ge$ the
calculated value.

\section{Alternative Design: Ferrite Core, Same Specification}
For comparison, the same \SI{800}{V}/\SI{800}{V}, \SI{10}{kHz}, \SI{250}{kW}
transformer is now sized with a MnZn ferrite core.

\subsection*{Assumptions}
\begin{itemize}
    \item MnZn ferrite (e.g.\ 3C95/3C90-type material), $B_m=\SI{0.3}{T}$
    (derated from $B_{sat}\approx\SI{0.4}{-}\SI{0.45}{T}$ at \SI{10}{kHz} and
    elevated temperature, to keep core losses bounded).
    \item Same $K_f=4$, $K_u=0.3$, $J=\SI{300}{A/cm^2}$ as before, for direct
    comparison.
\end{itemize}

\subsection*{Step 1 — Area product}
\begin{equation}
    A_p = \frac{500{,}000\times10^4}{4\times0.3\times0.3\times10{,}000\times300}
    \approx \SI{4630}{cm^4}
\end{equation}
This is roughly $3.3\times$ larger than the nanocrystalline design, a direct
consequence of the lower usable $B_m$ — illustrating why ferrite becomes
volumetrically penalized at high power and moderate (\SI{10}{kHz}) frequency,
despite its lower specific core loss.

\subsection*{Step 2 — Core area and turns}
\begin{equation}
    A_c \approx \sqrt{A_p} \approx \SI{68}{cm^2}
\end{equation}
\begin{equation}
    N_1 = \frac{800}{4\times68\times10^{-4}\times0.3\times10{,}000}
    \approx 9.8 \;\Rightarrow\; N_1=N_2=10 \ \text{turns}
\end{equation}

\subsection*{Step 3 — Conductors}
Current is unchanged ($I=\SI{312.5}{A}$, $A_{Cu}\approx\SI{1.04}{cm^2}$); skin
depth at \SI{10}{kHz} is also unchanged ($\delta\approx\SI{0.66}{mm}$), so the
same foil thickness limit ($\le\SI{1.3}{mm}$) applies, now over 10 turns instead
of 6.

\subsection*{Step 4 — Leakage inductance}
Window side $\approx\SI{8.25}{cm}\Rightarrow h_w\approx\SI{8}{cm}$;
$\mathrm{MLT}\approx\SI{30}{cm}$ for the larger core. Winding build-up with 10
foil turns:
\begin{equation}
    b_1=b_2=10\times(1.3+0.1)\,\text{mm}=\SI{14}{mm}
\end{equation}
With the same $\delta=\SI{2}{mm}$:
\begin{equation}
    L_{lk} \approx \frac{4\pi\times10^{-7}\times10^2\times0.30}{0.08}
    \left[\frac{0.014+0.014}{3}+0.002\right] \approx \SI{5.3}{\micro H}
\end{equation}
The higher turns count ($N^2$ scaling) and larger window both push leakage
inductance up by roughly $4.7\times$ relative to the nanocrystalline design —
a relevant trade-off if a specific $L_{lk}$ is targeted for soft-switching.

\subsection*{Comparison Summary}
\begin{center}
\begin{tabular}{@{}lll@{}}
\toprule
Parameter & Nanocrystalline & Ferrite (MnZn) \\
\midrule
$B_m$ & \SI{1.0}{T} & \SI{0.3}{T} \\
$A_p$ & $\approx\SI{1389}{cm^4}$ & $\approx\SI{4630}{cm^4}$ \\
$N_1=N_2$ & 6 & 10 \\
$L_{lk}$ (non-interleaved) & $\approx\SI{1.14}{\micro H}$ & $\approx\SI{5.3}{\micro H}$ \\
Relative core volume & baseline & $\approx 3.3\times$ \\
\bottomrule
\end{tabular}
\end{center}
At this power and frequency, nanocrystalline gives a markedly more compact
design with lower leakage; ferrite remains attractive mainly where its higher
achievable switching frequency (tens to hundreds of kHz) can be exploited, which
is not the case at \SI{10}{kHz}.

\begin{thebibliography}{9}

\bibitem{mclyman}
C.~W.~T.~McLyman, \textit{Transformer and Inductor Design Handbook}, 4th ed.,
CRC Press, 2011.

\bibitem{erickson}
R.~W.~Erickson and D.~Maksimovic, \textit{Fundamentals of Power Electronics},
3rd ed., Springer, 2020.

\bibitem{mohan}
N.~Mohan, T.~M.~Undeland, and W.~P.~Robbins, \textit{Power Electronics:
Converters, Applications, and Design}, 3rd ed., Wiley, 2003.

\bibitem{ferreira}
J.~A.~Ferreira, \textit{Electromagnetic Modelling of Power Electronic
Converters}, Kluwer Academic Publishers, 1989.

\bibitem{iec60664}
IEC~60664-1, \textit{Insulation coordination for equipment within low-voltage
systems -- Part 1: Principles, requirements and tests}, International
Electrotechnical Commission.

\bibitem{vacuumschmelze}
Vacuumschmelze, \textit{VITROPERM -- Nanocrystalline Tape-Wound Cores},
Application Notes and Material Datasheets.

\bibitem{ferroxcube}
Ferroxcube, \textit{Soft Ferrites and Accessories -- Data Handbook}, MnZn
Material Grades (3C90, 3C95, 3F3), Application Notes.

\bibitem{dowell}
P.~L.~Dowell, ``Effects of eddy currents in transformer windings,'' \textit{
Proceedings of the IEE}, vol.~113, no.~8, pp.~1387--1394, 1966.

\end{thebibliography}

\end{document}
